\section{Related Work}
\label{sec:related-work}

\subsection{Multi-hop QA and supporting-fact selection}

Multi-hop question answering evaluates the ability to combine evidence
across multiple documents. HotpotQA~\cite{yang2018hotpotqa} pairs each
question with ten Wikipedia paragraphs and sentence-level supporting-
fact annotations, making paragraph-selection quality directly
measurable. A line of work frames selection as a separate stage:
cognitive graph retrievers~\cite{ding2019cognitivegraph} and iterative
pathway retrievers~\cite{asai2020learning} traverse linked entities
across paragraphs, while transformer-based sentence
selectors~\cite{nie2019revealing} predict supporting sentences
directly. These methods are accurate but comparatively heavy. Our
TrustScore sits strictly upstream of any such selector: it is
pointwise, scores paragraphs in isolation, and does not reason over
the candidate graph.

\subsection{Retrieval-augmented generation and context scoring}

Retrieval-augmented generation~\cite{lewis2020rag,guu2020realm}
retrieves a small context set at inference time to ground a generator.
Dense Passage Retrieval~\cite{karpukhin2020dense} and downstream
variants such as Fusion-in-Decoder~\cite{izacard2021fid} and
REPLUG~\cite{shi2023replug} improved retrieval-side recall but
continued to score passages by a single learned relevance function.
Recent work has probed the failure modes of single-signal retrieval:
LLMs are sensitive to the position of the relevant
context~\cite{liu2024lostinthemiddle}, struggle with irrelevant
distractors~\cite{yoran2024makingrag}, and degrade further under
conflicting passages~\cite{xie2024adaptive}. Chen~et~al.\ introduced
the RGB benchmark~\cite{chen2024rgb} specifically to stress-test RAG
under noisy, counterfactual, and conflicting evidence. Our work
complements this line: rather than train a retriever end-to-end, we
ask whether a \emph{post-retrieval} trust-aware scorer, built from
cheap signals and a small learned combiner, can separate gold from
distractor candidates more reliably than a single relevance signal.

\subsection{Source attribution and trust in LLM outputs}

A parallel thread studies how language-model outputs should credit the
evidence they stand on. GopherCite~\cite{menick2022gophercite} and
Attributed QA~\cite{bohnet2022attributed} fine-tune generators to
produce quotation-backed answers; Gao~et~al.~\cite{gao2023enabling}
survey attribution in generative systems; Rashkin~et~al.~\cite{rashkin2023measuring}
define measurable attribution criteria; and Asai~et~al.'s
Self-RAG~\cite{asai2024selfrag} inserts explicit critique tokens that
assess evidence quality at decode time. These methods operate during
or after generation, and presume that a useful evidence pool has
already been selected. Our setting is upstream: we ask which
paragraphs deserve to enter the evidence pool in the first place, and
which signals predict that.

\subsection{Confidence, consistency, and hallucination filtering}

Trust-aware selection is closely related to confidence estimation and
hallucination detection. Self-consistency~\cite{wang2023selfconsistency}
uses agreement across sampled reasoning chains as a correctness
signal; selfcheckGPT~\cite{manakul2023selfcheckgpt} extends this to
hallucination detection; Kadavath~et~al.~\cite{kadavath2022language}
probe language models' internal calibration. These methods rely on
consistency across model samples. Our consistency signal is structurally
similar but operates across candidate \emph{passages} at retrieval
time, and our learned combiner discovers that in a distractor-rich
pool this signal is anti-correlated with gold, reversing the naive
reading of consistency as corroboration.

\subsection{Distinguishing contribution}

Most closely related is the broader idea, articulated in recent
position papers~\cite{wu2024retrieval,ram2023incontextrag}, that RAG
pipelines should treat sources as heterogeneous entities with
per-source trust rather than a homogeneous relevance pool. These
discussions are largely conceptual. We operationalize the idea on a
standard benchmark and show two concrete results. First, naive uniform
trust-signal averaging is \emph{actively harmful} on HotpotQA
distractor: not every trust heuristic is positively correlated with
gold, and averaging them without calibration is worse than using
relevance alone. Second, a six-feature logistic regression fit on 100
questions per seed recovers the correct sign of each signal, beats
SBERT by 3.3 absolute Precision@2 points, and more than doubles the
gold-versus-distractor score gap. The ablation isolates
a concrete instance of this phenomenon in HotpotQA distractor, and
we are not aware of an existing empirical demonstration at this
granularity.
